%=============================================================================!
% 11.1  HOLDING A BOND FIXED                                                  !
%=============================================================================!
\section{Holding a bond fixed}
\label{sec:cs-shake}

The fastest motions in most molecules are the stretches of bonds to
hydrogen (\cref{tab:os-vibrations}), and \cref{sec:pe-classical} found
them thoroughly quantum: for the O--H stretch of water $\hbar\omega$ is
17.5 times $\kB T$ at room temperature, so a simulation with classical
nuclei gets the energy of this vibration wrong whatever its step. Yet the
same vibration, with its period of \SI{9}{\femto\second}, is what holds
the step down. Replacing each such bond by a rod of fixed length
(\cref{sec:la-multipliers}) removes the vibration altogether. The bond
between atoms $i$ and $j$, held at the length $d$, is the constraint
\begin{equation*}
   \chi = \tfrac12\bigl(\lvert\vb{r}_{ij}\rvert^2 - d^2\bigr) = 0 ,
   \qquad \vb{r}_{ij} = \vb{r}_j - \vb{r}_i ,
\end{equation*}
whose force, found in \cref{sec:la-multipliers}, is $\lambda\vb{r}_{ij}$
on atom $j$ and $-\lambda\vb{r}_{ij}$ on atom $i$: along the bond, equal
and opposite, with a size set by the multiplier $\lambda$. This section
finds $\lambda$ at every step.

\subsection*{The constrained step}

The Störmer--Verlet step of \cref{eq:in-stormer}, written for atom $i$
with its force $\vb{F}_i$ and the constraint force $\vb{F}_i^{\mathrm c}$
added, is
\begin{equation*}
   \vb{r}_i(t + \dt) = 2\vb{r}_i(t) - \vb{r}_i(t - \dt)
   + \frac{\dt^2}{m_i}\bigl(\vb{F}_i + \vb{F}_i^{\mathrm c}\bigr) .
\end{equation*}
Without the constraint the atom would go to the provisional position
$\tilde{\vb{r}}_i = 2\vb{r}_i(t) - \vb{r}_i(t - \dt) +
(\dt^2/m_i)\vb{F}_i$, so the constrained step is the provisional one
corrected, $\vb{r}_i(t + \dt) = \tilde{\vb{r}}_i +
(\dt^2/m_i)\vb{F}_i^{\mathrm c}$. With one bond, and its force taken at
time $t$, the two atoms move to
\begin{align*}
   \vb{r}_j(t + \dt) &= \tilde{\vb{r}}_j
      + \frac{\dt^2\lambda}{m_j}\,\vb{r}_{ij} ,\\
   \vb{r}_i(t + \dt) &= \tilde{\vb{r}}_i
      - \frac{\dt^2\lambda}{m_i}\,\vb{r}_{ij} ,
\end{align*}
with $\vb{r}_{ij}$ the bond at time $t$. Subtracting the second line from
the first, the bond at $t + \dt$ is
\begin{align*}
   \vb{r}_{ij}(t + \dt)
   &= \tilde{\vb{r}}_{ij} + \dt^2\lambda\Bigl(\frac{1}{m_i}
      + \frac{1}{m_j}\Bigr)\vb{r}_{ij}\\
   &= \tilde{\vb{r}}_{ij} + \eta\,\vb{r}_{ij} ,
   \qquad \eta = \frac{\dt^2\lambda}{m_{\mathrm r}} ,
   && \text{$\frac1{m_i} + \frac1{m_j} = \frac{m_i + m_j}{m_im_j}
      = \frac{1}{m_{\mathrm r}}$,}
\end{align*}
where $\tilde{\vb{r}}_{ij} = \tilde{\vb{r}}_j - \tilde{\vb{r}}_i$ is the
provisional bond and $m_{\mathrm r}$ the reduced mass of
\cref{sec:os-bodies}. The new bond is the provisional one plus a multiple
of the old, as \cref{fig:cs-shake-geometry} draws. The constraint must
hold at $t + \dt$, $\lvert\tilde{\vb{r}}_{ij} + \eta\vb{r}_{ij}\rvert^2 =
d^2$. Expanding the square as $(\vb{a} + \vb{s})\cdot(\vb{a} + \vb{s})$
was expanded in \cref{sec:ro-freedom} gives $\lvert\tilde{\vb{r}}_{ij}
\rvert^2 + 2\eta\,\tilde{\vb{r}}_{ij}\cdot\vb{r}_{ij} + \eta^2\lvert
\vb{r}_{ij}\rvert^2 = d^2$, and with $\lvert\vb{r}_{ij}\rvert^2 = d^2$,
since the bond held at $t$,
\begin{equation}
   d^2\eta^2 + 2\,(\tilde{\vb{r}}_{ij}\cdot\vb{r}_{ij})\,\eta
   + \lvert\tilde{\vb{r}}_{ij}\rvert^2 - d^2 = 0 ,
   \label{eq:cs-quadratic}
\end{equation}
a quadratic equation for the one unknown $\eta$.

\begin{figure}[htb]
\centering
\begin{tikzpicture}[line cap=round, line join=round, scale=1.45,
   old/.style={circle, draw=black, fill=white, inner sep=1.8pt},
   trial/.style={circle, draw=black!45, fill=black!20, inner sep=1.8pt},
   new/.style={circle, fill=black, inner sep=1.8pt},
   arr/.style={-{Stealth[length=5pt]}, line width=0.8pt}]
   \draw[arr, black!60] (0,0) -- (3,0)
      node[midway, below] {$\vb{r}_{ij}$};
   \draw[black!45, dashed] (0.2,0.9) -- (4.0,1.5)
      node[pos=0.5, above=2pt, black!60] {$\tilde{\vb{r}}_{ij}$};
   \draw[line width=1pt] (0.8455,0.9) -- (3.7848,1.5);
   \draw[arr, accent] (0.2,0.9) -- (0.8455,0.9);
   \draw[arr, accent] (4.0,1.5) -- (3.7848,1.5);
   \node[old, label={below:$\vb{r}_i(t)$}] at (0,0) {};
   \node[old, label={below:$\vb{r}_j(t)$}] at (3,0) {};
   \node[trial, label={left:$\tilde{\vb{r}}_i$}] at (0.2,0.9) {};
   \node[trial, label={right:$\tilde{\vb{r}}_j$}] at (4.0,1.5) {};
   \node[new, label={below:$\vb{r}_i(t + \dt)$}] at (0.8455,0.9) {};
   \node[new, label={above:$\vb{r}_j(t + \dt)$}] at (3.7848,1.5) {};
\end{tikzpicture}
\caption{One step of a bond held at the length $d$, here 3. Atoms $i$
(mass 1) and $j$ (mass 3) start at the open circles. Without the
constraint they would reach the grey provisional positions, a bond
$\tilde{\vb{r}}_{ij}$ too long (dashed); the correction (arrows) moves
each along the old bond $\vb{r}_{ij}$, the lighter atom three times as
far, to the black positions, where the bond has length 3 again.}
\label{fig:cs-shake-geometry}
\end{figure}

\subsection*{One bond}

Two atoms, $i$ of mass 1 and $j$ of mass 3, joined by a bond of length 1
along $x$, drift apart to
1.2 in a step, so $\vb{r}_{ij} = (1, 0, 0)$, $\tilde{\vb{r}}_{ij} = (1.2,
0, 0)$, $\tilde{\vb{r}}_{ij}\cdot\vb{r}_{ij} = 1.2$ and
$\lvert\tilde{\vb{r}}_{ij}\rvert^2 - d^2 = 0.44$:
\begin{align*}
   \eta^2 + 2.4\,\eta + 0.44 &= 0 ,\\
   \eta &= \frac{-2.4 \pm \sqrt{2.4^2 - 4\times0.44}}{2}
   = -1.2 \pm 1 ,
   && \text{the quadratic formula,}
\end{align*}
so $\eta = -0.2$ or $-2.2$. The first shortens the bond to $1.2 - 0.2 =
1$. The second gives $1.2 - 2.2 = -1$, the bond turned round, a solution
of the equation that no short step can reach: as $\dt$ shrinks,
$\tilde{\vb{r}}_{ij}$ approaches $\vb{r}_{ij}$ and the correction must
vanish, which only the root nearer zero does: with $\tilde{\vb{r}}_{ij} =
\vb{r}_{ij}$, \cref{eq:cs-quadratic} becomes $d^2\eta^2 + 2d^2\eta = 0$,
with roots 0 and $-2$. With $\eta = -0.2$ and
$m_{\mathrm r} = 1\times3/(1 + 3) = 0.75$, atom $j$ moves by
$\dt^2\lambda/m_j = \eta m_{\mathrm r}/m_j = -0.05$ along the bond and atom
$i$ by $+0.15$: the lighter atom moves three times as far, and the centre
of mass, $1\times0.15 + 3\times(-0.05) = 0$, does not move at all.

The quadratic has real roots only if the number under the square root of
the quadratic formula, $4(\tilde{\vb{r}}_{ij}\cdot\vb{r}_{ij})^2 -
4d^2(\lvert\tilde{\vb{r}}_{ij}\rvert^2 - d^2)$, is not negative, that is,
if $(\tilde{\vb{r}}_{ij}\cdot\vb{r}_{ij})^2 \ge d^2(\lvert\tilde{\vb{r}}_{ij}
\rvert^2 - d^2)$. Writing
$\phi$ for the angle between the provisional and the old bond, so that
$\tilde{\vb{r}}_{ij}\cdot\vb{r}_{ij} = \lvert\tilde{\vb{r}}_{ij}\rvert
\,d\cos\phi$, and dividing by $d^2$, the condition reads
$\lvert\tilde{\vb{r}}_{ij}\rvert^2\cos^2\phi \ge \lvert\tilde{\vb{r}}_{ij}
\rvert^2 - d^2$, that is, $d^2 \ge \lvert\tilde{\vb{r}}_{ij}\rvert^2(1 -
\cos^2\phi) = \lvert\tilde{\vb{r}}_{ij}\rvert^2\sin^2\phi$, or, taking
square roots, $\lvert\tilde{\vb{r}}_{ij}\rvert\sin\phi \le d$. The part of the provisional bond at
right angles to the old one must be no longer than the bond: a correction
along the old bond cannot shorten that part. A bond stretched by 50\% and
turned in one step by more than $\arcsin(1/1.5) = \ang{41.8}$, and less
than $\ang{180} - \ang{41.8} = \ang{138.2}$, where $\sin\phi$ is again
$1/1.5$, cannot be corrected at all, a sign of a step far too long.

\subsection*{Many bonds}

The bonds of a molecule share atoms. The three distances held in a rigid
water molecule join its three atoms in a triangle, and moving two atoms to
correct one distance changes the other two, so the quadratics of the three
bonds are coupled. Ryckaert, Ciccotti and Berendsen avoided solving them
together\cite{Ryckaert1977}: their method, SHAKE, corrects the bonds one
at a time, each by an approximate solution of its own equation with the
others' atoms where they are, and sweeps through all the bonds again and
again until every one holds. The approximation drops the $\eta^2$ of
\cref{eq:cs-quadratic}, which is small for a short step:
\begin{equation}
   \eta \approx \frac{d^2 - \lvert\tilde{\vb{r}}_{ij}\rvert^2}
   {2\,\tilde{\vb{r}}_{ij}\cdot\vb{r}_{ij}} .
   \label{eq:cs-linear}
\end{equation}
For the example this gives $-0.44/2.4 = -0.1833$, and the bond becomes
1.0167 long. Applied again, with the corrected bond in place of
$\tilde{\vb{r}}_{ij}$ and $\vb{r}_{ij}$ still the old bond, it gives
$-0.0165$, and
the bond is $1.000137$. The repeated step is a method for solving any
equation.

\begin{toolbox}[Newton's method]
To solve $f(y) = 0$, we start from a guess $y_n$ and replace the curve
$f(y)$ near it by its tangent, the first two terms of its Taylor series
(\cref{sec:mo-taylor}), $f(y_n) + f'(y_n)(y - y_n)$. The tangent is zero at
\begin{equation*}
   y_{n+1} = y_n - \frac{f(y_n)}{f'(y_n)} ,
\end{equation*}
the next guess. To see how fast the guesses close in, we write the root
as $y^*$ and the error as $e_n = y_n - y^*$, and expand $f(y^*)$ about
$y_n$, with $y^* - y_n = -e_n$:
\begin{align*}
   0 = f(y^*) &= f(y_n) - e_nf'(y_n) + \tfrac12e_n^2f''(y_n)
   + \order{e_n^3} ,\\
   \frac{f(y_n)}{f'(y_n)} &= e_n - \frac{f''(y_n)}{2f'(y_n)}\,e_n^2
   + \order{e_n^3} ,
\end{align*}
the second line by dividing the first by $f'(y_n)$ and rearranging.
Subtracting $y^*$ from both sides of the rule gives $e_{n+1} = e_n -
f(y_n)/f'(y_n)$, and so
\begin{equation*}
   e_{n+1} = \frac{f''(y_n)}{2f'(y_n)}\,e_n^2 + \order{e_n^3} .
\end{equation*}
Each error is about a fixed multiple of the square of the one before, so
once the error is small the number of correct digits roughly doubles at
every step. For $f(y) = y^2 - 2$ from $y_0 = 1$, the guesses $1.5$,
$1.4166667$, $1.4142157$ and $1.4142135624$ have errors of $0.086$,
$0.0025$, $\num{2.1e-6}$ and $\num{1.6e-12}$ against $\sqrt2$. The method
needs $f'(y_n) \neq 0$ and a start close enough to the root; far from it,
the tangent can point anywhere.
\end{toolbox}

For one bond, $f(\eta) = \lvert\tilde{\vb{r}}_{ij} +
\eta\vb{r}_{ij}\rvert^2 - d^2$, the left side of \cref{eq:cs-quadratic},
so $f'(\eta) = 2d^2\eta + 2\,\tilde{\vb{r}}_{ij}\cdot\vb{r}_{ij} =
2\,\vb{r}_{ij}\cdot(\tilde{\vb{r}}_{ij} + \eta\vb{r}_{ij})$,
which at $\eta = 0$ is $2\,\tilde{\vb{r}}_{ij}\cdot\vb{r}_{ij}$, so
\cref{eq:cs-linear} is the first step of Newton's method from $\eta = 0$.
Once the atoms have moved, the corrected bond takes the place of
$\tilde{\vb{r}}_{ij}$, and the same formula is the next step. Here
$f''(\eta) = 2d^2$ and, near the root, $f'(\eta)$ is twice the dot product
of the old bond with the corrected one, $2d^2\cos\theta$ if the two bonds
are $\theta$ apart, so each error in $\eta$ is $2d^2/(2\times2d^2\cos
\theta)$ times the square of the one before. An error $e$ in $\eta$ changes
the length by about $e\,d\cos\theta$, so the relative error of the length
is $1/(2\cos^2\theta)$ times the square of the one before: about half,
when the bond has not turned far. \Cref{fig:cs-shake}(a) measures it: a
bond stretched by 10\% has a relative error of 0.1, then $\num{5.0e-3}$,
$\num{1.4e-5}$ and $\num{1.1e-10}$, each 0.50 to 0.57 times the square of
the one before; the corrected bond lies \ang{19.9} off the old one, and
$1/(2\cos^2\ang{19.9}) = 0.565$.

\begin{figure}[htb]
\centering
\includegraphics{ch11_constraints/shake}
\caption{How SHAKE converges. (a) One bond of length 1 between equal
masses, stretched by 1\%, 10\% and 30\% in a direction \ang{18} off the
old bond: the relative error of its length after each sweep, until it
reaches the rounding of the arithmetic. (b) The 64 rigid water molecules
of \cref{sec:cs-water} after one step of \SI{2}{\femto\second} without
constraints: the sweeps needed against the tolerance (circles) and the
largest error of a held distance that is left (squares, right axis).}
\label{fig:cs-shake}
\end{figure}

The function \texttt{shake} of \texttt{mdlab.constraints} carries out the
sweeps. The bonds are sorted beforehand into groups that share no atom, so
that every bond of a group can be corrected at once with NumPy. For
water, whose bonds are listed molecule by molecule, this applies the same
corrections in the same order as taking the bonds one at a time; for other
lists two bonds that share an atom can change places within a sweep, which
changes the path of the sweeps but not the conditions they end on. The
loop is
\begin{code}
    for sweep in range(1, max_sweeps + 1):
        moved = False
        for g in groups:
            i, j = pairs[g, 0], pairs[g, 1]
            now = r[j] - r[i] + copy[g]
            eta = 0.5 * (d2[g] - np.sum(now * now, 1)) / np.sum(
                old[g] * now, 1
            )
            eta[np.abs(eta) <= tolerance] = 0.0
            moved = moved or bool(np.any(eta))
            push = (m_r[g] * eta)[:, None] * old[g]
            r[j] += push / m[j, None]
            r[i] -= push / m[i, None]
        if not moved:
            return r, sweep
\end{code}
\noindent Here \texttt{old} holds the bonds $\vb{r}_{ij}$ at time $t$,
\texttt{now} the current ones, \texttt{d2} the squared lengths and
\texttt{m\_r} the reduced masses; \texttt{push} is $m_{\mathrm r}\eta\,
\vb{r}_{ij}$, which divided by $m_j$ is atom $j$'s move. The array
\texttt{copy} holds, for each bond, the multiple of the lattice vectors
that turns $\vb{r}_j - \vb{r}_i$ into its nearest copy in the provisional
positions (\cref{sec:md-image}), so a molecule that lies across a face of the cell
is corrected as one piece. A sweep that leaves every $\lvert\eta\rvert$
below the tolerance moves nothing, and ends the loop. Since $\eta$ is
about the fraction by which a bond is too long or too short (a bond $1 +
s$ times its length along the old one gives, by \cref{eq:cs-linear}, $\eta
= -s(1 + \frac12s)/(1 + s) \approx -s$), the
tolerance is a relative error of the lengths: $10^{-10}$ in
\texttt{mdlab} unless another is asked for.

Several bonds converge more slowly than one. \Cref{fig:cs-shake}(b) takes
the box of 64 rigid water molecules of \cref{sec:cs-water} through one
step of \SI{2}{\femto\second} with no constraint, which leaves the
largest held distance in error by \SI{0.030}{\angstrom}, 3.1\% of an O--H
bond, and corrects it by SHAKE. A tolerance of $10^{-10}$ needs 39 sweeps.
Each tenfold reduction of the tolerance costs 4.5 more sweeps, so each
sweep reduces the error by a steady factor $q$ with $q^{4.5} = 10$, $q =
10^{1/4.5} = 1.67$, where a
single bond squared it: correcting one side of a triangle disturbs the
other two, and the errors are passed round it. The largest error left is
1.5 times the tolerance in \si{\angstrom}, as it should be for a relative
error of distances up to the H--H length of \SI{1.51}{\angstrom}.

\begin{keyidea}
A held bond adds a force along the old bond to the step, so the new bond
is the provisional one plus $\eta$ times the old. The constraint at the
end of the step is a quadratic in $\eta$, and its first-order solution is
a step of Newton's method. SHAKE takes the bonds one at a time with that
step, and sweeps until every bond holds to a tolerance.
\end{keyidea}

\notebook{11}{drift a bond and follow SHAKE sweep by sweep}

\begin{selfcheck}
In the example with masses 1 and 3, the masses are changed to 2 and 2.
Does $\eta$ change, and how far does each atom move?
\end{selfcheck}
